\documentclass[10pt,conference]{IEEEtran}
\usepackage[T1]{fontenc}
\usepackage{microtype}
\usepackage{amsmath,amssymb,booktabs,array,tabularx}
\usepackage[hyphens]{url}
\usepackage[hidelinks]{hyperref}
\usepackage{enumitem}
\setlist{nosep,leftmargin=*}
\newcommand{\code}[1]{\texttt{#1}}
\hypersetup{pdftitle={From Model Metadata to Executable Interfaces: A Controlled Study of Tensor Binding and Decoder Contracts}}
\title{From Model Metadata to Executable Interfaces:\\A Controlled Study of Tensor Binding and Decoder Contracts}
\author{\IEEEauthorblockN{Anonymous author(s)}}
% Workshop working draft; venue choice, author audit and artifact delivery pending.
\begin{document}
\maketitle


\begin{abstract}
Successful inference does not establish every property of a generated model interface. We characterize this boundary through controlled associated-label interventions on an EfficientNet classifier. Deleting one label preserves graph/weights, embedded metadata, generated Java/classes and tested raw outputs, yet the actual Android category getter throws after inference. A count-preserving permutation instead succeeds while changing 999 of 1,000 returned strings; score bits remain fixed. Reversal restores behavior. On the same Android virtual device, an official Task classifier rejects the identical deletion at construction. Existing decoder guards and a public generated-getter failure precede this work. Our contribution is controlled public-interface evidence distinguishing structural acceptance, reference-backed association and selected-decoder compatibility, with preservation checked separately. The results identify integration-test boundaries; they establish neither a new general checker nor defect prevalence.
\end{abstract}

\section{Question and Contract Boundaries}
\label{sec:contracts}
A generated model interface consumes more than numerical tensors: it can load associated labels and expose them alongside scores. TensorFlow Lite Support documents metadata-driven interface generation~\cite{schema,codegen-doc}. We ask: when computation is preserved, how can these resources change an executed public interface, and where do configured official consumers enforce the relevant condition?

The known constraint is not the contribution. A 2022 public report, Support \#770, already describes a segmentation getter reaching a TensorLabel mismatch. A repository-member reply calls that model's generation route unsupported and recommends Task Library~\cite{support770,support770reply}. We supply controlled classifier executions with a working original, negative control, preservation and reversal, followed by a same-Android comparison. The reported segmentation artifact is an antecedent, not a natural failure reproduced here.

Write the package as $(G,W,O,M,F)$: graph, weights, ordered outputs, embedded metadata and associated files. We distinguish:
\begin{itemize}
\item \textbf{S: structural acceptance} by a specified entry point, here Support 0.4.4 MetadataPopulator loading and repopulating metadata on a disposable copy. Reading and populating are separate operations.
\item \textbf{I: reference-backed association agreement}, comparing output roles or display-label order with an independently sourced, artifact-linked specification. Without one, I abstains; shape or the original labels alone are insufficient.
\item \textbf{D: selected-decoder compatibility}, concerning the buffer, resource and overload actually used by the getter. For Support 0.4.4's list-based TensorLabel category path, the first axis with size greater than one must exist, have size equal to the loaded label count, and be the final axis~\cite{tensorlabel}.
\end{itemize}
These are different questions, not certification levels. \textbf{P}, separately, checks intervention preservation: computational fields and non-metadata buffers, canonical graph after metadata removal, embedded metadata, associated ZIP payloads, generated code and observed outputs. Whole-package hashes should change under a resource intervention.

\section{E1: An Executed Getter Failure}
\label{sec:e1}
The focal artifact is the version-1 float32 EfficientNet Lite0~\cite{efficientnet}, SHA-256 prefix \code{6c7ab0a6e5dc}, with input \code{[1,224,224,3]}, output \code{[1,1000]} and 1,000 associated labels. Four conditions are original; one-character description control; deletion of the last label; and restoration of the label payload using the same deterministic ZIP writer. Deletion changes the appended associated file while preserving embedded metadata. Restoration recovers the payload, not the publisher's original ZIP-container bytes.

Support 0.4.4 generates Java that compiles unchanged. The APK uses JDK 17 targeting Java 8, TFLite 2.13.0, Support 0.4.4 and an explicit build harness, rather than the generated experimental/nightly Gradle project unchanged. Execution uses Android API 30, x86\_64, CPU and one thread. The harness calls
\begin{quote}\small\ttfamily\raggedright
ContractModel.newInstance(...)\\
\(\rightarrow\) process(TensorImage)\\
\(\rightarrow\) Outputs.getProbabilityAsCategoryList()
\end{quote}
It observes a duplicate output buffer before the actual getter without replacing the decoder. Two inputs were fixed before inspection: RGB 127 and row-major bytes $(37j+17)\bmod256$, each $224\times224\times3$, passed through UINT8 TensorImage and generated default preprocessing. The unit is component $\times$ variant $\times$ consumer/configuration; inputs are repeated observations.

All E1 conditions pass S, generation, source compilation, packaging, loading and inference. Original, description and restoration return 1,000 categories for both inputs. Deletion alone throws at the generated getter:
\begin{quote}\small\raggedright
\code{IllegalArgumentException}: ``Label number 999 mismatch the shape on axis 1''.
\end{quote}
The stack enters the generated Outputs getter; the harness catches it. This is a failed public call, not a demonstrated application crash or a natural EfficientNet defect.

All 165 non-metadata buffers and canonical graph agree. Deletion/restoration preserve metadata and the pre-ZIP prefix. Generated Java/classes and corresponding raw-output bytes agree across conditions; complete APKs differ because their model assets differ. These controls locate the consequence after preserved numerical computation (Table~\ref{tab:core}).

\begin{table}[t]
\caption{Five core model-byte conditions, both fixed inputs. Graph/weights, generated Java/classes and raw outputs agree. Description alone changes embedded metadata. Shared phase aliases are not independent models.}
\label{tab:core}
\centering\footnotesize
\begin{tabular}{@{}lrll@{}}
\toprule
Condition & Labels & Getter & Exposed strings\\
\midrule
Original & 1,000 & success & baseline\\
Description & 1,000 & success & same\\
Deletion & 999 & exception & not returned\\
Permutation & 1,000 & success & 999 changed\\
Restoration & 1,000 & success & restored\\
\bottomrule
\end{tabular}
\end{table}

\section{E2: Successful but Changed Association}
\label{sec:e2}
E2 predeclares the cyclic order $[L_0,\ldots,L_{C-1}]\mapsto[L_1,\ldots,L_{C-1},L_0]$, retaining filename, count and multiset. Original, description, permutation and inverse restoration all complete the real getter. Raw bytes, score bits, argmax and top-five indices agree; 999/1,000 returned strings change under permutation. Duplicate adjacent ``maillot'' strings explain the one unchanged position. Inverse restoration recovers the original lists.

For the pattern input, output position 904 retains score \code{0.23542051017284393}, while ``window screen'' becomes ``window shade''. Positions are verified against the raw tensor; this overload leaves the Category index field unset. Computation is preserved while the exposed association changes.

\paragraph{Qualified independent reference.}
\label{sec:p3}
A separate publisher-reference comparison connects the focal computation to TFHub Lite0 fp32/2: ordered computational fields and all 165 non-metadata buffers agree despite different packages. Its pinned catalog card identifies parent classification/1, whose instructions prescribe an external label file with the first background entry removed~\cite{tfhub-lite-card,tfhub-parent-card,imagenet-display}. All 1,000 original/restoration positions match; permutation disagrees at 999. This supports $I_{\rm display}$ disagreement while S and D accept. The catalog is a later retrieved snapshot and the external label URL is unversioned but byte-pinned. Stable class-ID/training lineage remains unresolved; no prediction-accuracy conclusion follows.

Primary E1/E2 comprises eight phase aliases of five byte conditions, one graph and 16 input observations (14 lists, two caught failures). A separate agent construction audit regenerated/rebuilt in new directories and a new AVD using installed tools, reproducing outcomes and matching generated sources/classes. Recorded replay agrees; an interrupted replay attempt remains separate. These are not independent human reproduction.

\section{Same-Android Enforcement Stages}
\label{sec:consumers}
P2 runs the same five model-byte conditions in both consumers, through separate wrapper and Task ImageClassifier APKs on the same API-30 x86\_64 AVD, CPU/one thread. The wrapper retains Support 0.4.4/TFLite 2.13.0; Task Vision/Base 0.4.4 uses its published native integration. All ten configurations build and install. Task deletion fails at construction before either input, leaving 18 actual input observations.

\begin{table}[t]
\caption{Same-Android public boundaries. The four successful Task constructions also complete classify on both inputs. Deletion never reaches Task classify.}
\label{tab:stages}
\centering\footnotesize
\begin{tabular}{@{}llll@{}}
\toprule
Condition & Wrapper process & Wrapper getter & Task create\\
\midrule
Deletion & pass & fail & fail\\
Other four & pass & pass & pass\\
\bottomrule
\end{tabular}
\end{table}

Task's outer IllegalStateException reports an error obtaining a native address; its IllegalArgumentException cause identifies 1,000 classes versus 999 labels. Versioned source locates that comparison in initialization's \code{CheckAndSetOutputs}~\cite{task-classifier}. The generated route instead constructs TensorLabel at getter access. Both harnesses catch exceptions. Across eight successful shared condition/input pairs, all returned scores/labels agree after aligning Task class indices to wrapper output positions.

This controls device, OS, API and ABI while retaining native-integration differences and an unknown embedded Task interpreter version. Task classify is atomic at our boundary; its raw tensor is not separately captured. The result locates enforcement in two configured paths, not an isolated interpreter effect or a measured safety/latency benefit.

\section{Boundary Cases and Supporting Checks}
\label{sec:hand}
Hand provides a secondary limit. Metadata names output 0 \code{handedness} at \code{[1,63]} and output 2 \code{landmarks} at \code{[1,1]}, opposite the roles in pinned official integration source. The source applies through a separate source-pinned component with equal ordered computation and 172 non-metadata buffers~\cite{hand-graph,hand-assets}. Swapping entries changes generic bindings, but both generated sources require the same disclosed separator repair before compilation/disassembly. Official MediaPipe 0.10.14 Task decoded fields remain identical across original, container control, swap and restoration on one fixed image. This is no second Android-wrapper execution or universal insensitivity result; the swapped scalar still violates the selected generic decoder's nonunit-axis requirement.

\label{sec:preflight}
An external purposive examination froze criteria/checker and 20 URLs before inspection: four unavailable, 16 downloads, 18 extracted occurrences and three prior component/graph exclusions leave 15 new graphs. All pass S and abstain under I; eight wrappers compile, seven fail, and none is executed. D reports four compatible, nine incompatible and two inapplicable paths; the nine already have compilation failures or unsupported-output warnings. They are not nine natural runtime defects.

The template-specific E5 preflight adds zero rejections over a matched simple decoder baseline or Task initialization on the five core conditions: each rejects deletion only. Reading and population accept all five; a separate output-metadata count reduction is rejected by population. Dynamic-signature abstentions include a Gesture deletion with a known isolated-probe failure. A later reconstruction report also exposes opposite Python/Java label-parsing errors: a vertical-tab mutation is accepted by the Python approximation but rejected by Java, while a U+2028 mutation produces the reverse. These are isolated decoder/initialization probes, not Android inference; their raw records are currently unrecovered. Thus parser fidelity, unsupported code and evidence scope constrain this supporting diagnostic; no general analyzer or detector-accuracy contribution is claimed.

\section{Implications, Prior Work and Validity}
\label{sec:discussion}
The practical result is test placement. A trusted whole-package hash would detect these resource changes, but cannot judge an intentionally approved replacement. Raw regression observes computation; an actual getter test exposes E1; reference-map comparison exposes E2's display disagreement. Task construction already enforces cardinality. We measure coverage boundaries, not operational savings or prevalence. E1/E2 concern one classifier and two synthetic inputs; mutants and replay do not expand the independent-model count.

\label{sec:related}
Existing work covers ML contracts/wrappers (MLGuard), documentation-derived constraints (DocTer), implementation-derived constraints (ACETest), model/naming consistency (DARA), mutation testing (DeepMutation) and numerical differential testing (NNSmith)~\cite{mlguard,docter,acetest,dara,deepmutation,nnsmith}. Our narrower empirical pairing holds computation and generated code fixed while observing actual output-access failure or changed display association, then compares configured enforcement stages. Neither mutation nor the decoder rule is new; \#770 precludes a first-getter-failure claim.

\label{sec:validity}
\subsection{Data Availability and Author Verification}
Source and claim indexes are retained. A verified local package contains 2,749 selected original E1/E2, construction-audit and P2/P3 payload files; saved checks cover 26 executed APKs and 34 complete records including replay, not extra independent models. Recovery does not rederive graph preservation, compiled bindings or reference authority. Hand, corpus and checker/baseline raw families are outside this focal package; later reconstruction/parser records remain unrecovered. Persistent anonymous evidence access is pending, so this working draft is not submission-ready. No model experiment was rerun for this revision. AI assisted code, execution, analysis and writing; human author verification and final disclosure remain pending.

\section{Conclusion}
Preserved computation coexists with either a post-inference getter exception or successful but changed display associations. Reversal and the same-Android comparison locate these outcomes at concrete consumer operations. The evidence motivates observing the selected public interface alongside numerical outputs, while keeping structural acceptance, reference agreement and decoder compatibility distinct.

\begin{thebibliography}{18}\small\raggedright
\bibitem{schema} TensorFlow Lite Support. \emph{Metadata schema}, v0.4.4. \url{https://github.com/tensorflow/tflite-support/blob/v0.4.4/tensorflow_lite_support/metadata/metadata_schema.fbs}.
\bibitem{codegen-doc} Google AI Edge. \emph{Generate model interfaces with metadata}. Retrieved October 2026; snapshot preserved. \url{https://developers.google.com/edge/litert/android/metadata/codegen}.
\bibitem{tensorlabel} TensorFlow Lite Support. \emph{TensorLabel.java}, v0.4.4. Selected source/binary identities are recorded in the companion index. \url{https://github.com/tensorflow/tflite-support/blob/v0.4.4/tensorflow_lite_support/java/src/java/org/tensorflow/lite/support/label/TensorLabel.java}.
\bibitem{efficientnet} MediaPipe Models. \emph{EfficientNet Lite0, float32, version 1}. \url{https://storage.googleapis.com/mediapipe-models/image_classifier/efficientnet_lite0/float32/1/efficientnet_lite0.tflite}.
\bibitem{tfhub-lite-card} TensorFlow Hub. \emph{EfficientNet Lite0 fp32/2 model card}. Catalog commit \code{81fdea8ffa53b5acbd66a012d17fb8f9602c714b}. \url{https://github.com/tensorflow/tfhub.dev/blob/81fdea8ffa53b5acbd66a012d17fb8f9602c714b/assets/docs/tensorflow/models/efficientnet/lite0/fp32/lite/2.md}.
\bibitem{tfhub-parent-card} TensorFlow Hub. \emph{EfficientNet Lite0 classification/1 model card}. Same catalog commit; explicit output-index mapping and background rule. \url{https://github.com/tensorflow/tfhub.dev/blob/81fdea8ffa53b5acbd66a012d17fb8f9602c714b/assets/docs/tensorflow/models/efficientnet/lite0/classification/1.md}.
\bibitem{imagenet-display} TensorFlow. \emph{ImageNetLabels.txt}. Retrieved reference snapshot, SHA-256 prefix \code{536feacc519d}; full bytes, digest, and storage generation preserved. \url{https://storage.googleapis.com/download.tensorflow.org/data/ImageNetLabels.txt}.
\bibitem{task-classifier} TensorFlow Lite Support. \emph{ImageClassifier Java/JNI/C++ implementation}, v0.4.4, commit \code{133228616ca14e1b64c03f6a0ea8610f41a60783}. Companion source/dependency identities record the selected integration. \url{https://github.com/tensorflow/tflite-support/blob/133228616ca14e1b64c03f6a0ea8610f41a60783/tensorflow_lite_support/cc/task/vision/image_classifier.cc}.
\bibitem{hand-graph} MediaPipe. \emph{Hand landmarks detector graph}, v0.10.14, resolved commit \code{4cf89a70942c}. Immutable-commit source identified in the historical record. \url{https://github.com/google-ai-edge/mediapipe/blob/4cf89a70942ca3252e46ace7e4552f53be9bef2e/mediapipe/tasks/cc/vision/hand_landmarker/hand_landmarks_detector_graph.cc}.
\bibitem{hand-assets} MediaPipe. \emph{External file declarations}, commit \code{4cf89a70942ca3252e46ace7e4552f53be9bef2e}. Source-pinned model/image identities recorded in the historical evidence index. \url{https://github.com/google-ai-edge/mediapipe/blob/4cf89a70942ca3252e46ace7e4552f53be9bef2e/third_party/external_files.bzl}.
\bibitem{mlguard} Sheng Wong et al. MLGuard: Defend Your Machine Learning Model! In \emph{Proceedings of the 1st International Workshop on Dependability and Trustworthiness of Safety-Critical Systems with Machine Learned Components (SE4SafeML)}, 2023. \href{https://doi.org/10.1145/3617574.3617859}{doi:10.1145/3617574.3617859}.

\bibitem{docter} Danning Xie et al. DocTer: Documentation-Guided Fuzzing for Testing Deep Learning API Functions. In \emph{Proceedings of the 31st ACM SIGSOFT International Symposium on Software Testing and Analysis (ISSTA)}, 2022. \href{https://doi.org/10.1145/3533767.3534220}{doi:10.1145/3533767.3534220}.

\bibitem{acetest} Jingyi Shi et al. ACETest: Automated Constraint Extraction for Testing Deep Learning Operators. In \emph{Proceedings of the 32nd ACM SIGSOFT International Symposium on Software Testing and Analysis (ISSTA)}, 2023. \href{https://doi.org/10.1145/3597926.3598088}{doi:10.1145/3597926.3598088}.

\bibitem{dara} W. Jiang et al. ``I see models being a whole other thing'': an empirical study of pre-trained model naming conventions and a tool for enhancing naming consistency. \emph{Empirical Software Engineering}, 30:155, 2025. \href{https://doi.org/10.1007/s10664-025-10711-4}{doi:10.1007/s10664-025-10711-4}.
\bibitem{deepmutation} Lei Ma et al.
DeepMutation: Mutation Testing of Deep Learning Systems.
\href{https://arxiv.org/abs/1805.05206v2}{arXiv:1805.05206v2}, 2018.

\bibitem{nnsmith} Jiawei Liu et al.
NNSmith: Generating Diverse and Valid Test Cases for Deep Learning Compilers.
In \emph{Proceedings of the 28th ACM International Conference on Architectural
Support for Programming Languages and Operating Systems (ASPLOS)}, Volume 2,
2023.
\href{https://doi.org/10.1145/3575693.3575707}{doi:10.1145/3575693.3575707}.
Inspected version: \href{https://arxiv.org/abs/2207.13066v2}{arXiv:2207.13066v2}.

\bibitem{support770} TensorFlow Lite Support issue \#770.
Semantic Segementation Example ``Label number 21 mismatch'' exception.
Opened March 28, 2022; issue and comments retrieved October 4, 2026.
\url{https://github.com/tensorflow/tflite-support/issues/770}.

\bibitem{support770reply} lu-wang-g.
Reply to TensorFlow Lite Support issue \#770, April 25, 2022.
GitHub author association: MEMBER; retrieved October 4, 2026.
\url{https://github.com/tensorflow/tflite-support/issues/770#issuecomment-1108777779}.
\end{thebibliography}
\end{document}
